\section{Cơ sở lý thuyết và khoảng trống}

\begin{frame}{Từ rigid job đến moldable job}
\begin{columns}[T,onlytextwidth]
\begin{column}{0.31\textwidth}
\begin{block}{Rigid}
Một cấu hình duy nhất:
\[
p_j=p_j^{original}
\]
Scheduler chỉ quyết định thời điểm và vị trí chạy.
\end{block}
\end{column}
\begin{column}{0.31\textwidth}
\begin{block}{Moldable}
Nhiều cấu hình khả thi:
\[
\mathcal{P}_j=\{p_{j1},\ldots,p_{jk}\}
\]
Chọn profile trước khi job bắt đầu, allocation sau đó giữ nguyên.
\end{block}
\end{column}
\begin{column}{0.31\textwidth}
\begin{block}{Malleable}
Tài nguyên thay đổi trong runtime:
\[
p_j=p_j(t)
\]
\end{block}
\end{column}
\end{columns}
\vspace{0.4em}
Nghiên cứu này tập trung vào \textbf{moldability trước submission}, không thay đổi allocation khi job đang chạy.
\end{frame}

\begin{frame}{Moldable scheduling đã chứng minh giá trị của resource flexibility}
Các nghiên cứu thường giải bài toán:
\[
\mathcal{P}_j\;\longrightarrow\;\text{profile selection}\;\longrightarrow\;\text{scheduling}
\]
\begin{itemize}
    \item lựa chọn số processor hoặc node cho job;
    \item đồng thời chọn profile và thời điểm bắt đầu;
    \item tối ưu makespan, waiting time hoặc turnaround time;
    \item multi-resource, QoS-aware và preference-aware scheduling;
    \item online scheduling cho moldable jobs và workflows.
\end{itemize}
\begin{alertblock}{Kết luận từ literature}
Nhiều resource alternatives có thể cải thiện makespan, utilization và throughput khi scheduler được phép nhìn thấy chúng.
\end{alertblock}
\end{frame}

\begin{frame}{Giả định phổ biến: profile set đã được biết}
Các phương pháp moldable scheduling thường nhận:
\[
\mathcal{P}_j=\{p_{j1},p_{j2},\ldots,p_{jk}\}
\qquad\text{hoặc}\qquad
t_j(p)
\]
\begin{itemize}
    \item profile do application developer khai báo;
    \item performance curve có được từ profiling;
    \item đủ historical executions;
    \item hoặc profile space đã được định nghĩa trong mô hình toán học.
\end{itemize}
\vspace{0.3em}
\textbf{Câu hỏi còn thiếu:} Làm sao hình thành các alternatives đáng tin cậy từ một request ban đầu vốn rigid?
\end{frame}

\begin{frame}{Các nhánh related work chưa được nối liền}
\begin{enumerate}
    \item \textbf{HPC scheduling:} tối ưu thứ tự, placement và backfilling nhưng thường nhận request đã có cấu trúc.
    \item \textbf{Moldable scheduling:} khai thác nhiều profile nhưng giả định profile set đã tồn tại.
    \item \textbf{Prediction và profiling:} dự đoán runtime, memory hoặc utilization nhưng khó với job mới và ít dữ liệu.
    \item \textbf{Natural-language intent:} chuyển ngôn ngữ tự nhiên thành cấu trúc nhưng chưa chứng minh feasibility.
    \item \textbf{LLM-assisted scheduling:} có reasoning đa mục tiêu nhưng không tự bảo đảm consistency và policy compliance.
\end{enumerate}
\end{frame}

\begin{frame}{Research gap: tạo ra profile có thể kiểm chứng trước scheduler}
\begin{block}{Known}
Moldable scheduling có thể tối ưu hệ thống khi tập resource alternatives đã được cung cấp.
\end{block}
\begin{block}{Limitation}
Các phương pháp hiện tại giả định đầu vào đã có cấu trúc, profile set hoặc performance evidence đủ mạnh.
\end{block}
\begin{alertblock}{Unknown}
Chưa rõ làm thế nào chuyển
\[
\text{rigid request}+\text{incomplete intent}+\text{cluster constraints}+\text{limited evidence}
\]
thành tập profiles vừa khả thi kỹ thuật, phù hợp intent, được cho phép và hữu ích cho scheduler.
\end{alertblock}
\end{frame}
